% ============================================================
%  Harald Høffding, RECOLLECTIONS [Erindringer]
%  Copenhagen: Gyldendalske Boghandel, Nordisk Forlag, 1928.
%  TRANSLATION (English) OF THE TRANSCRIBED SECTIONS ONLY.
%
%  Translated from transcription.tex (Danish) in this directory, whose
%  header records how its text was established from the 1928 first
%  printing (every page read at the image by transcriber and second
%  reader, with a blind control).  Like the transcription, this file is a
%  SKELETON of the whole book: only pp. 39--51, 55--69, 77--82, 90--95
%  and the table of contents are translated; every other stretch stands
%  as a \gap line and a "% [text to be added: pp. X--Y]" marker, to be
%  filled when the Danish is transcribed.  Method:
%  ../../../TRANSLATION-PLAYBOOK.md.  The earlier English made from the
%  2021 e-book (translation.tex.bak.ebook-withdrawn-20260923) was not
%  opened; no published translation or other edition was consulted.
%
%  ---- REGISTER ----
%  Readable modern English, moderately literal; American spelling.  The
%  book is a memoir, written 1915--24, and is rendered as one; the letters
%  of 1866--70 quoted on pp. 63--69 are rendered as letters.  Long periods
%  are broken only where English will not carry them, never across a page
%  marker or a footnote anchor.
%
%  ---- TERMINOLOGY ----
%  Aligned with the other translations in this repo:
%    Tro / Viden = faith / knowledge;  Erkendelse = cognition, but
%    Erkendelsesteori = epistemology;  Videnskab = science
%    (Wissenschaft), videnskabelig = scientific;  Aand = spirit;
%    Kristendom = Christianity, Urkristendom = primitive Christianity;
%    Livsanskuelse = view of life;  den Enkelte = the single individual;
%    Privatdocent = Privatdozent;  Vurderingsmotivet / Handlingsmotivet
%    = the motive of valuation / the motive of action.
%  Kierkegaard's works in their standard English titles, inside the
%  author's quotation marks ("Stadier" = ``Stages'', Kierkegaard's
%  „Øjeblikke“ = issues of ``The Moment'', „Til Selvprøvelse“ = ``For
%  Self-Examination'').  All other titles stay in their original language
%  as printed, including the titles of lecture courses.  Names unchanged;
%  kings as Frederik VII, Karl XV; Copenhagen, Zealand, North Zealand,
%  Funen, Jutland, Gothenburg in English, other place names as printed.
%  Hr. / Fru = Mr. / Mrs.; Danish titles of rank untranslated.  The dates
%  of the letters are given with the month names in full.
%
%  ---- CONVENTIONS ----
%  - Page markers \opage{N} and "% --- p. N ---" at the same joints as in
%    transcription.tex; a word the Danish divides across a page turn is
%    kept whole in English, with the marker before or after it.
%  - The same paragraphs, \textit spans (the 1928 book's emphasis is
%    italic), footnotes (\footnote[N] with the printed number), white
%    lines (\whiteline) and rules (\secrule) as the Danish.
%  - Quotation marks „...“ -> ``...''.  Latin, German and French stay in
%    their language; Greek verbatim.
%  - Where a translated range begins or ends inside a sentence whose other
%    part lies on an untranscribed page, the English begins or ends with
%    an editorial \ldots.
%  - Translator's decisions are recorded in "% TR:" comments at the site.
%
%  ---- DEFECTS: DANISH AS PRINTED, ENGLISH CORRECT ----
%  The transcription is diplomatic and carries every printer's error.  The
%  translation renders the evident SENSE and logs the divergence in a
%  % TR: comment: p. 41 "andelig" (= aandelig), "Kierkegaard „Til
%  Selvprøvelse“" (genitive); p. 44 "sprængtes.Hinsidighedens"; p. 46
%  "mine legemlige Konstitution"; p. 47 "ligened"; p. 59 "Rasmus
%  Nielsens hørte"; p. 60 "dennne", "( som"; p. 63 "ngen"; p. 78 "sine
%  unge ... Hustru", "Tillidt"; p. 79 "Hovedtænker", rendered "principal
%  ideas" (Hovedtanker; the print, confirmed at the image, has æ);
%  p. 92 a clause without a subject ("he" supplied).  p. 61 "tænkte" is
%  partly supplied by the editor where the scan is lost.
%
%  ---- HOW THE TRANSLATION WAS CHECKED (2026-09-23) ----
%  (1) Six translators, one per range (pp. 39--45, 46--51, 55--60,
%  61--69, 77--82, 90--95), each re-reading its English against the
%  Danish sentence by sentence before returning.  (2) Mechanical audit
%  (tr_audit.py, adapted from ../../brochner/erindringer-kierkegaard/):
%  per printed page the same \opage joints, \textit/\emph spans,
%  footnotes, white lines, rules, chapter heads and paragraph breaks as
%  transcription.tex, and an English/Danish word ratio within
%  0.95--1.60: 40 pages, 0 flagged.  (3) Independent review: three
%  readers who had not translated the pages read every sentence of the
%  Danish against the English.  No omission of substance found; 29
%  corrections applied, among them p. 39 a dropped "here" and modal;
%  p. 46 "maatte" (was bound to, not could); p. 62 the greatest
%  systematic thinker of the century is Comte, not Hegel, and the
%  pluperfect "havde gjort"; p. 64 a relative clause mistaken for a
%  concession; p. 79 "Hovedtænker" (above); p. 81 "gerne" = usually.
%
% ============================================================
\documentclass[11pt]{book}
\usepackage[utf8]{inputenc}
\usepackage[T1]{fontenc}
\usepackage{libertinus}
\usepackage{libertinust1math}
\usepackage{textalpha}
\usepackage{amsmath}
\usepackage[english]{babel}
\usepackage{setspace}
\usepackage{fancyhdr}
\setlength{\headheight}{28pt}
\addtolength{\topmargin}{-13.5pt}
\pagestyle{fancy}
\fancyhf{}
\fancyhead[LE,RO]{\thepage}
\fancyhead[RE]{\textit{Harald Høffding: Recollections}}
\fancyhead[LO]{\textit{\leftmark}}
\setstretch{1.3}

% Footnotes as the 1928 printing marks them: 1), 2), ... numbered afresh
% on each printed page, written \footnote[N]{...} with the printed number.
\renewcommand{\thefootnote}{\arabic{footnote})}

% Editorial: a stretch of the book not yet transcribed or translated.
\newcommand{\gap}[1]{\par\begin{center}\small[#1 not yet transcribed or translated]\end{center}\par}
% A white line in the print, and a short centred rule.
\newcommand{\whiteline}{\par\vspace{\baselineskip}\par}
\newcommand{\secrule}{\par\begin{center}\rule{2cm}{0.4pt}\end{center}\par}

\usepackage{hyperref}
\hypersetup{
  pdftitle={Recollections},
  pdfauthor={Harald Høffding},
  hidelinks
}

\usepackage{marginnote}
\newcommand{\opage}[1]{\marginnote{\footnotesize\textit{[#1]}}}

\begin{document}

\frontmatter
\begin{titlepage}
\centering
\vspace*{2cm}
{\small\emph{HARALD HØFFDING}}\par
\vspace{1.5cm}
{\LARGE\emph{RECOLLECTIONS}}\par
\vspace{0.5cm}
{\small [\emph{ERINDRINGER}]}\par
\vspace{1cm}
{\large\textasteriskcentered}\par
\vfill
\rule{0.6\textwidth}{1.2pt}\par\vspace{1pt}\rule{0.6\textwidth}{0.3pt}\par
\medskip
{\small\emph{GYLDENDALSKE BOGHANDEL - NORDISK}}\par
{\small\emph{FORLAG - KØBENHAVN - MCMXXVIII}}\par
\bigskip
{\small Translated from the Danish (transcribed sections only)}\par
\end{titlepage}

\chapter*{Preface}
\addcontentsline{toc}{chapter}{Preface}
\gap{Preface (p.~[6])}
% [text to be added: pp. 6--6]

\mainmatter

% ---- Chapter I.  Indhold: "I. 1843—1861 .... 7".  pp. 7--37 (PDF 9--39).
\chapter*{I}
\addcontentsline{toc}{chapter}{I. 1843--1861}
\markboth{I. 1843--1861}{I. 1843--1861}
\gap{Chapter I, pp.~7--37,}
% [text to be added: pp. 7--37]

% ---- Chapter II.  Indhold: "II. 1861—1874 .... 38".  pp. 38--93 (PDF 40--95).
\chapter*{II}
\addcontentsline{toc}{chapter}{II. 1861--1874}
\markboth{II. 1861--1874}{II. 1861--1874}
\gap{p.~38}
% [text to be added: pp. 38--38]
% p. 39 opens mid-sentence: "de bestaa sammen!" ("they exist together!"); p. 38 not yet transcribed
% --- p. 39 ---
\opage{39}\ldots{} they exist together! Must it not, in the end, be possible here
to find one single expression for the life of the personality? ---

It was Rasmus Nielsen and Sibbern who became my first
teachers in philosophy. Rasmus Nielsen was then at the
height of his powers. He was still searching, captivated by Kierkegaard
and absorbed in mathematical and natural-scientific
studies, striving to find a way by which these different
worlds of thought could be united. His propaedeutic
lectures in my year consisted partly of a survey of
% TR: „Stadier“ (and „Stadierne“ below) is the author's short form of Stages on Life's Way; rendered ``Stages''.
Kierkegaard's ``Stages,'' on the basis of a small selection he
had had printed, and partly of a brief epistemological outline.
He spoke with life and force, often brilliantly and wittily,
and I owe those lectures an extraordinary amount.
They set a host of thoughts in motion in me which
had never stirred before. With that course of lectures, however, it went with him
as it so often went with his books: the plan
was too large and comprehensive, or else he worked out
its first parts at such length that there was not
time enough left for the later ones. Thus of the ``Stages'' he
got through only the first, the ``aesthetic.'' At every
possible opportunity he polemicized against theology as a science,
sharply distinguishing between faith and theology.
At first this polemic roused my opposition.
I had come with great expectations of what the study of
theology would bring; I hoped that the isolation
in which the religious had so far stood in the life of my soul
would, with the help of theology, give way to a harmony among
the different spheres. I hoped to be led into
a union of the life of thought and the life of feeling, such that they would complement
% TR: „Atene“ = the goddess Athene (paired with the Madonna), not the city.
each other; I hoped that Athene and the Madonna could
be united into one. In a diary I kept in my first years as a student
I accordingly spoke harshly of those who would
not acknowledge that there is a Christian science. --- That
was soon to change.

While Rasmus Nielsen lectured on philosophical propaedeutics,
Sibbern lectured on psychology and logic. (The philosophical
% --- p. 40 ---
\opage{40}course then took up ten hours a week. That
was perhaps too much. But they have gone too far in
limiting it to four hours, even if, with the use of
printed textbooks, not a little can now be accomplished in that time. In
my opinion six hours would be suitable.) If propaedeutics
was kept separate from psychology and logic, it was because it
% TR: Erkendelsesteori = epistemology (term of art), though Erkendelse is elsewhere "cognition"; almindelig Videnskabslære = general theory of science.
was properly meant to provide epistemology, or a general theory of science.
Rasmus Nielsen, incidentally, gave a good deal besides
epistemology under the heading of propaedeutics; thus,
in my year, the survey of the Kierkegaardian stages.
There was a third professor, Hans Brøchner; he gave courses parallel to
Sibbern's and Nielsen's by turns, and in my year he lectured on
Sibbern's subjects. I had chosen Sibbern, because I wanted to
hear him before he retired, and because I expected a great deal
of the author of ``Gabrielis.'' He was then 77 years old.
It proved a disappointment. There were, to be sure, flashes of genius now and then,
beautiful thoughts and priceless remarks; but his
delivery was flat and mechanical, and his dulled senses
let him neither hear how he himself rattled his
keys while he talked, nor notice the restlessness and chatter
of his audience. His naivety and oddity could be priceless,
as when he named his own ``Gabrielis'' among
``good poetic works,'' or when he explained to us at length
why the four clocks of Frue Kirke did not agree. But when,
because he was too warm up at the lectern, he first took
off his skullcap, and, when that did not help, his wig, one could
not ask young people to take it with reverence.
He had outlived himself as a teacher, although his inner
life still stirred with freshness. His books were heavy going
to learn, all the more since he liked best to have his own words given back to him.
--- Only much later, when I had myself become independent in
philosophy, did I properly learn to appreciate him, and it has
been a great joy to me to bring his memory forward and to characterize
him\footnote[1]{See ``Mindre Arbejder,'' first series, and ``Danske Filosofer.''}, and to help make his thoughts
known abroad as well.

% --- p. 41 ---
\opage{41}In my first year as a student I also began my theological
studies, especially exegesis and church history. The exegetical
study is probably what I have profited from most
within theology. It accustomed one to analyzing what one
read, to exact and thorough understanding, and it gave
occasion for interesting comparisons among the New Testament
writers. And to become at home in so
great a work of the spirit as the New Testament, to be able to wander
through it crisscross and to feel the connection of the thoughts
with the linguistic form, was already a great
gain, and it is this above all that has made me never
regret having studied theology, however far I have
come away from theological ways of thinking. To be sure, the lengthy
German commentaries we studied were no delight
to get sense out of. One plowed through a whole series of
interpretations which in the author's eyes were wrong, and
then at last reached the author's own reading, which of course
was the height of wisdom. There was a pedantry in this
that could give one a distaste for the study. And when one really stopped to
consider that the words which were the object of so much
learning were meant to set those who read them face to face with
life's greatest question, the question of spiritual
% TR: printed „andelig“ (for aandelig); rendered by its evident sense.
life or spiritual death, then one could often be indignant at
the whole manner of treatment, especially when one had just
% TR: printed „Kierkegaard“ without the genitive s (for Kierkegaards); rendered by its evident sense.
read a book such as Kierkegaard's ``For Self-Examination.'' More than
once I flung the commentaries across the floor --- but I
went over afterwards and picked them up. After all, I meant to finish what
I had begun. ---

Through all this, however, I lived for the time being as a cheerful
student and had many pleasant evenings in the old
Students' Association building in Boldhusgade, for all its squalor.
% TR: „Sold paa Salen“ = a student drinking-party in the hall; rendered ``revels in the hall.''
On Saturday evenings there were often ``revels in the hall.''
One sat with one's pipe by the punch bowl; the songbook with
Ploug's and Hostrup's songs was opened, and between the songs
more or less stirring or witty speeches were made.
% TR: Senioratet = the Seniorate, the elected governing board of the Students' Association.
Discussions were not held, except at general meetings,
% --- p. 42 ---
\opage{42}where battles were fought to topple or uphold the Seniorate.
In the mornings too one could expect to meet
comrades up there, and I got much use out of the good
library.

Since there were not as many students as now, and
since they had more lectures in common, there was ample opportunity
for making acquaintances. I quickly came to
know the students from Aalborg School, several of whom became
good friends of mine. Not many days had passed after
lectures began before I was on intimate terms
with Andreas Faartoft\footnote[1]{Later district physician in Nykøbing on Mors.}, quite a type of the young student
of that time, open and cheerful and full of zest for life.
He had a fine singing voice, which often gave pleasure in
my home, where he soon became a welcome guest. Another
man from Aalborg, Søren Svendsen\footnote[2]{Later principal of the navigation school on Bogø.}, I also quickly got to
know, and he and his violoncello likewise soon became
welcome guests in my home. My parents and brothers and sisters
regarded my two young comrades as belonging
at home with us whenever there was to be a pleasant and
cozy evening. I owe them much, not least in the
following years, when I was no enlivening element
in my home, and when I myself nevertheless needed enlivening.
Faartoft later became a respected physician in his native town, Nykøbing
on Mors, and Svendsen, who had studied physics, became
head of the navigation school on Bogø. --- A third man from Aalborg,
Poul Kierkegaard\footnote[3]{Son of Bishop P. Chr. Kierkegaard.}, I came into a closer
relationship with somewhat later, and I shall tell in another connection of
him and his fate. ---

% TR: Filosofikum = the examen philosophicum, the preliminary philosophy examination taken by all students.
On the same day that I had sat for the Filosofikum,
Norwegian and Swedish students arrived. It was the last Scandinavian
student meeting to be held in Copenhagen (1862).
There was great enthusiasm in the city and among the students,
and the days that followed were glad and pleasant ones, days I
% --- p. 43 ---
\opage{43}still think of as festive memories, although dark shadows soon
fell over the hopes that formed
the ground of the festive mood. I especially remember
the excursion in North Zealand, when, having spent the night in Hillerød,
we walked the next morning along the glorious road from Hillerød
to Fredensborg. Here there was opportunity for conversation and
exchange of thoughts with our Scandinavian brothers, and one
could learn about conditions of study at the other Nordic
universities. Already on the first day I was quite taken
with a young Swede whose name was Svartling, and
with whom I corresponded for several years. Many years later
I heard him described as a dean of some standing. At Fredensborg
we filed through the domed hall of the palace, from whose
gallery Frederik VII, surrounded by Countess Danner,
ministers and courtiers, received our homage. We marched
% TR: „Sexa“ (Swedish sexa) = a light cold meal; rendered "cold collation".
down into Nordmandsdalen, where a cold collation had been laid out.
But hardly had we begun to help ourselves to the dishes when
a violent thunderstorm broke loose. We ducked down under
the tables, taking with us whatever we could lay hands on of food and
drink, especially the latter. No wonder the mood
was high when the sun came out again and we emerged once more.
A little later the King came, and in his genial way
he said he hoped we had got as wet inside
as outside. The excursion then went on to Helsingør.

That summer my family lived in a house on Krogerup
Allé, from which there was a glorious view across the fields to
the Sound and to Helsingborg. It is one of the best summer holidays
% TR: in note 1, p. 43, classical names in their English forms (Aristophanes, Socrates, Plato); titles as printed.
I have had\footnote[1]{I see from my letters to Hans Skovgaard from the summer holidays of 1862
that I occupied myself a great deal with Aristophanes (in Danish translation),
particularly in order to get clear about the relation between him and Socrates.
My interest had already turned on Socrates since my school days. Besides this I
read Plato's ``Gorgias'' (partly in Greek, partly in Danish), Rötscher's ``Aristophanes
und sein Zeitalter,'' Goethe's ``Aus meinem Leben,'' and Kierkegaard's
religious discourses. It was the combination of these occupations with
excursions in that glorious countryside that made the holidays in Humlebæk unforgettable.}. We roamed about the glorious countryside,
now on horseback and now on foot. The tenant farmer at Krogerup
% --- p. 44 ---
\opage{44}had two Icelandic horses, which he let my younger brother
and me have, and so we could get far and wide in
pleasant fashion. --- One day we drove to Danstrup Hegn
near Fredensborg and there saw Frederik VII, at the head
of his Horse Guards, receive Karl XV of Sweden,
who came driving from Helsingør and now mounted his horse,
whereupon both kings, at the head of a brilliant retinue,
rode to Fredensborg. This meeting of kings was thought, like
the student meeting before it, to herald a new
age for the North. Alas, they were soap bubbles, and heavy
and dark years came for our youthful ideals.

\whiteline

Now a heavy and dark time came in my inner
life as well. I immersed myself in Kierkegaard's writings, and after
a hard struggle I had to grant that he was right in his conception
of Christianity. It no longer stood as the
mild, consoling way of truth that could take up into itself all
human ways of life --- family life, art, science,
social life --- and lead them on toward perfection. The idyll
% TR: printed „sprængtes.Hinsidighedens“ with no space after the full stop; rendered as two sentences.
was shattered. The ideal of the beyond came to stand in all its
unconditional severity. Only one thing, in the literal
sense of the word, was needful: to enter into relation with the
% TR: Menneskeblivelsens Under = the miracle of God's becoming man (the Incarnation); kept in the author's non-technical phrasing.
great paradox, the miracle of God's becoming man, which made all that is
human pale and set only one task, that of giving up
everything in order to share in the divine powers which,
at one single place in the world and at one single point in time,
had broken through the continuity of nature and of human life
and let the light of eternity shine into the darkness. Sayings
in the New Testament that I had until then pushed aside
or interpreted mildly and harmonizingly now came to stand in
their whole and absolute ideality. It was not the intention of Christianity
to complete what human culture had begun,
but to break off, and to set aside all human
relations as indifferent, insofar as they were not hindrances
to reaching the great goal. It was as if the scales fell
from my eyes when I now took up the New Testament again.
% --- p. 45 ---
\opage{45}And it stood clear to me that what was proclaimed there was the
highest. I cannot point to any definite moment of conversion.
It was after hard resistance that I gave up
the harmonious and idyllic faith of my early youth. That
Christianity was, and would be, consolation, the eternal consolation,
was true; but this consolation was bound to one condition:
full devotion to the one thing needful, full renunciation of
the human purposes of life. Only now, in fact, did I discover
that neither art nor science nor civic activity
is mentioned in the New Testament, and that there are warnings
in stern words against marriage. It was Christian ethics
that now appeared to me in an entirely new light. With dogmatics
there was, for the time being, nothing amiss.

How, then, was life to be led, now that the ideal had thus
shown itself to lie higher than I had thought? The
time that now came did bring me moments of heartfelt enthusiasm,
but above all it brought me doubt, unrest and homelessness.
I was at my wits' end. The master, the prophet,
had after all passed away, after he had said that the Christianity
of the New Testament did not exist, if indeed it had ever
existed beyond the very earliest time. It was vain
to speculate on what position he would have taken
had he not been carried off so early. And besides,
he himself had taught me that each individual must solve
the problem of life in his own way --- that at bottom one cannot
learn from others.

And now my theological studies! I worked zealously at
the various theological subjects, but also continued my philosophical
studies, kept on attending Rasmus Nielsen's evening lectures,
and now also Brøchner's on the history of philosophy.
% TR: Privatissimum = a private course for a few students; kept as the author's (unitalicized) Latin term.
I took part in a privatissimum that Brøchner held on
Kant's Kritik der Urtheilskraft. But from the purely philosophical
side it was maintained ever more firmly that a scientific harmony
between faith and knowledge was impossible --- that at bottom
there was no theological science. So I wandered
back and forth between theology and philosophy --- also
% p. 46 opens mid-sentence, not mid-word: p. 45 ends "ogsaa" ("also"), p. 46 begins
% "bogstaveligt talt" ("literally"); the paragraph continues from p. 45 (no indent)
% p. 46 continues mid-sentence from p. 45 ("--- also | literally speaking"); the paragraph continues (no indent)
% --- p. 46 ---
\opage{46}literally speaking, since they had their lecture halls at opposite
ends of the long corridor upstairs in the University. A constant
intellectual discord made itself felt in me, besides the personal
one called forth by the new conception of Christianity.

Those were the heaviest and darkest years of my life. I shut myself
up with my doubts and became unsociable and closed toward those
nearest to me. However much I cared for them, I was not a good
member of the family. Life in my room and life in the living room
came to stand before me as two worlds, and I managed only all too
poorly to make the passage between them. The bright and harmonious
way in which those downstairs took everything irritated me. Precisely
in those years greater prosperity came to the house, and people
enjoyed, in their way, whatever it could now open the way to. I
held myself aloof and cold. I have many sins on my conscience from
these years, and only later have I understood how much they must
have had to bear with me. My parents, as I have since learned, were
worried about me and asked my friends for advice.

Through all this there also worked a tendency to hypochondria,
% TR: printed "mine legemlige Konstitution" (printer's error for "min"); rendered by the sense, "my bodily constitution".
which is bound up with my bodily constitution and has followed me
all my life, right up into old age. Only those nearest to me have
had occasion to notice this tendency, which with strangers I as a
rule find it easy to hold back. In my book ``Den store Humor'' the
section that deals with hypochondria (acedia) is one of the parts
most built on my own experience. Heaviness and darkness could,
without any demonstrable cause, spread over the mind and give
everything a dull and misty look. For a long time a feeling would
often come over me as though I were hanging by a hair that might
snap at any moment. Dissatisfied with myself, and at accomplishing
nothing in the world, I often felt that everything was hopeless for
me. This tendency was naturally bound to gain the upper hand all the more
easily as the mind stood doubting and uneasy before the greatest
% --- p. 47 ---
\opage{47}questions. How the balance of strength lay between the two
currents, the one that came from the bodily heaviness and the other
that came from the weight of the intellectual and the personal
problems, I naturally cannot untangle, and could do so even less
then.

But naturally it was not the case that I went about constantly like
a lost soul. My studies occupied me a great deal. Besides my own
subject, the interest in Hebrew that had already been awakened at
school led me to read Arabic with Professor van Mehren. I got so far
% TR: "Tusind og een Nat" given its standard English title (a translated work, not a Danish book title).
as to read a couple of tales from the ``Thousand and One Nights'' in
the original language, and I even thought for a while of becoming an
Orientalist, which would presumably have drawn me away from my
absorption in philosophical and religious problems. I still read a
good deal of philosophy; in particular I followed the development
from Kant to Hegel, a period that stood out with a peculiar luster
especially at that time. For not only had Kierkegaard, after all,
been strongly influenced by Hegel, but both Rasmus Nielsen and
Brøchner represented shades of Hegelianism. I also read some of
Plato's dialogues in the original language. In short, my time was
well occupied.

Systematic theology I read together with Hans Skovgaard, as a rule
out at his home on Valby Bakke. Our reading often stretched far
into the night, and I then had a walk before me in pitch darkness
from the hill down through the narrow lane that led from Langgade
to Rahbeks
% TR: printed "ligened" (closed up); read as "lige ned", "straight down".
Allé, and straight down to Vesterbrogade, where civilization began
with its street lamps. The road was not without danger, since
assaults often took place outside the city, and since at the country
house just opposite Pilealleen (now the main building of Ny
Carlsberg) some large dogs ran loose, which rushed out at the
nocturnal wanderer without his being able to see them. But it was
a refreshing walk I had out there in the afternoon. I remember a
summer afternoon when I came as usual to go through an agreed
section. I met Skovgaard's second-youngest sister in the garden
room, and she --- always
% --- p. 48 ---
\opage{48}amiable --- showed me a quite special friendliness that day
and drew me into a long and lively conversation. Afterwards it came
to light that she was acting by agreement with her brother, who had
not finished his assigned reading. She often had to take teasing for
the comedy she had played. She was a healthy and cheerful young
girl, and I often think of her when I walk past the old Bjerregaard.
Some years ago I had her granddaughter, who greatly resembled her,
among my listeners at the University. --- It was both enjoyable and
instructive to read with Hans Skovgaard. He was very interested,
but at the same time very skeptical and contrary. At that time I
was really a more eager theologian than he; later it has become the
other way round. In particular he was not as eager as I to find
meaning in Professor I. A. Bornemann's often obscure words. In my
experience, reading together with fellow students is something that
cannot be recommended enough to young students. It develops
independent thinking and criticism more than reading with a
% TR: "Manuduktør" (a private coach for examinations) rendered "tutor".
tutor, which, when the tutor is not especially gifted as a teacher,
easily takes on a mechanical character.

Besides Valby there was another place to which I often walked out,
and where I found rest and understanding. An old uncle of mine lived
in a country house in Frederiksberg Smallegade, and his house was
kept by a distant relative of my parents, Sofie Pape, sister of the
young girl whom I had visited toward the end of my journey in Jutland\footnote[1]{The two sisters were daughters of the landowner H. C. Pape of the
farm Heiselt in Vendsyssel. For a time he also held a lease near Viborg.
(See above p. 36).}.
In her I found an older sister, such as I had often wished for. She
understood my interest in Kierkegaard, though she kept mainly to his
religious discourses. And now, in the winter of 1863---64, her
younger sister Emma came to visit her. The three of us came to form
a little clique. But it was soon for the younger one's sake that I
went there. We read together and we walked together. And I now
thought
% --- p. 49 ---
\opage{49}that I had gained a sister of my own age. That it was a
feeling of another kind that was stirring in me became clear to me
only after her departure. For the time being I pushed back
everything that stirred in that direction, under the influence of
the strict religious discipline to which I had submitted myself. And
yet, far more than I would admit to myself, the thought of her was
in those dark years a gleam of light that kept courage and hope up.
Here, moreover, lay the seed of something that later worked to bring
me into a living relation to life again.

The winter of '63---64 was dark and heavy for my country. People
followed the events with suspense as they unfolded after the death
of Frederik VII: the threatening Federal Execution, the
revolt in Holstein, and the intervention of Prussia and Austria.
Then in February the battles at Dannevirke and the army's sudden
retreat. Grief and indignation seized everyone, and after the fall
of Dybbøl the defeat stood as inevitable. They were bitter and dark
times. But I am ashamed to think how much, through all this, I was
nevertheless absorbed in my inner struggles and broodings and let
myself be ruled by the feeling of powerlessness that followed from
them, without the misfortune of the country being able to rouse me
to action. In the early summer, however, I decided, together with my
friend Skovgaard, to report for enlistment. For some time, then, it
stood before me that now everything was to be laid aside, and it
also became clear to me that if I came back from the war, I would
not return to theology. --- Since no more soldiers were being
accepted, however, I returned to my study.

I had not given up Christianity, but I made to myself the admission
that Kierkegaard (in ``Practice in Christianity'') had demanded of
the representatives of the Church. I acknowledged that I did not
% TR: "Urkristendom", "urkristelig" rendered "primitive Christianity", "of primitive Christianity" throughout.
live under the ethics that primitive Christianity presupposes or
entails. I recognized my distance from the ideal of primitive
Christianity, although I maintained that it was the ideal. I felt
like a Catholic
% --- p. 50 ---
\opage{50}who does not enter a monastery, although he admits that it
is the ideal. I thought that on these terms I had the right to
regard myself as a Christian, although I could not think of becoming
a pastor.

The Church as Church could not --- I have since understood this
--- make the admission Kierkegaard demanded. That would have been
the same as admitting that it represented an entirely different
religion from that of primitive Christianity, once a religion and
its ethics are recognized at all as belonging so closely together
that they stand and fall with each other. But the single individual
who makes the admission also, after all, really acknowledges ---
this too I did not see clearly then --- that he does not have the
religion the Church claimed to represent. In reality it signified a
break with the Church by virtue of one's own experience of life and
thinking, and I thereby placed myself under the principle of
personality. The scope of this dawned on me only late and slowly, in
connection with my continued experience of life and my continued
studies. On this point Kierkegaard himself, by virtue of the
principle that subjectivity is the truth, points far beyond himself.
My love for Kierkegaard and my admiration for him have therefore been
able to stay with me long after I had moved away from all dogmas.
Yet it has been said of me in German, apropos of my book on
Kierkegaard, by a lightweight theologian, that I lack love for
Kierkegaard and immersion in him. But the fact is that there are
people who do not understand a love and an admiration that vent
themselves not in sentimental expressions, but in taking up a
problem that has been set and working at it faithfully in the way
that accords with the single individual's personality. Kierkegaard's
problem has followed me from my youth, determined the direction of
my life, and again and again led me to a more inward conception of
manifold relations, to a stricter testing of myself, to a shrinking
from whatever lacks personal truth. And while theologians and
half-theologians, in the centenary of his birth, whitewashed the
prophet's tomb, I held
% --- p. 51 ---
\opage{51}to the great principle of subjectivity as truth and sought
to apply it in the way that accorded with my personality and my
studies.

Purely philosophically, the study of Kierkegaard led me to a thought
which my later philosophical study led me to deepen and apply in
detail: the thought, namely, that the formal peculiarity of the life
of the spirit lies in the unity, the concentration, with which the
content of life is gathered together, and that the measure of the
worth of a spiritual life lies in the relation between the fullness
of the content and the energy of the concentration. This thought,
which had emerged in the history of philosophy before Kierkegaard,
underlies Kierkegaard's philosophical characterization of the
various stages or types of life. This point, naturally, was not
clear to me either. It dawned on me only through a critical study
of English philosophy and through the preliminary work for my
Psychology.

There was another standpoint from which in those years I could not
yet judge Kierkegaard's conception of primitive Christianity either.
Theologians are accustomed to regard the thought that carries the
whole ethics of primitive Christianity, namely the expectation of
the imminent coming of the ``Kingdom of Heaven,'' as an inessential
illusion of the earliest congregation, or they explain away the
passages where it is mentioned with particular force. And to my
great astonishment I saw, when Kierkegaard's posthumous papers were
printed, that he himself regarded the expectation as an illusion ---
only he reversed the relation: in his opinion one could not live by
an ethics such as that of primitive Christianity without that
illusion: ``It is to such a degree agonizing to be the true
Christian that it would be unbearable if one did not constantly
expect Christ's return right now.'' (Efterladte Papirer. 1849,
p. 248). By calling the expectation an illusion Kierkegaard places
himself in complete opposition to primitive Christianity. From a
historical standpoint the expectation is a fact and is the constant
positive presupposition, \ldots
% p. 51 ends mid-sentence at "Forudsætning," (presupposition,); p. 52 not yet transcribed
\gap{pp.~52--54}
% [text to be added: pp. 52--54]
% p. 55 opens mid-word: Danish "Ved-/kommende" (p. 54 ends "hans Ved-"); p. 54 not yet transcribed
% --- p. 55 ---
% TR: the slice opens on the second half of "Vedkommende" ("[for] hans Vedkommende", in his case); rendered "case" and preceded by \ldots{} because p. 54 is not yet transcribed.
\opage{55}\ldots{} case. He set a bundle of cigars on his table and
decided that they should be the last he smoked. This
last pleasure he then put off day after day, and in the course
of a year the craving had disappeared.---

The examination period went by as if it were sport, and on 19 June 1865
I walked down the University steps as candidatus theologiæ.
In the afternoon I had the conversation with my father
about my future that I mentioned earlier. My
wish was to go into school teaching and to continue
my studies. Even before I had gone
out to my family, who were staying in the country at Skodsborg, I received an offer
% TR: school names kept as printed (``Mariboes Skole'' here, ``Maribos Skole'' p. 56; the Danish spells it both ways).
of lessons at ``Mariboes Skole'' (by Frederiksholms Kanal)
in history and Danish, which I had particularly had in mind
as my teaching subjects. At this school I worked
as a teacher for 17 years, right up until I became professor,
and I have many good memories from that time. For the time being
I took a holiday at Skodsborg, and later in the summer
I made a journey to Stockholm, from where I returned home
% TR: "Gøteborg" (the Danish form of Göteborg) rendered "Gothenburg".
by the canal route by way of Gothenburg.

They were hard problems that I took with me from the
period of my life that had now closed. I had kept
my personal faith, with the reservation I have described before.
But in what relation should, and could, this
faith now stand to whatever else occupied me in life? That was
the main question that followed me. And to answer
it, philosophical thinking was necessary. Philosophy, which
had until then been a subsidiary subject for me, now became a main subject.
For the answer to that question required an investigation
of what human cognition is and what it is capable of, and in what
relation it stands to personal life. I had to
learn what the human life of the spirit is able to reach on
its own, so as then to be able to see whether a connection
between it and religious faith was possible, and on what
terms. That was the problem the Kierkegaard studies had
set before me, but had not solved. I had to pursue
this question further. My future I pictured to myself
% --- p. 56 ---
\opage{56}for the time being thus: that teaching would be my
outward work, and that to it would be joined the studies
it required, --- but that I would then strive to keep an
% TR: "Lønkammer" (the "closet" of Matt. 6:6) rendered "inner chamber".
inner chamber where what lay close to my heart personally
could be cultivated. That my philosophizing could come to have any
significance whatsoever for others did not occur to me.

And so, after the summer holidays, I began my work as a teacher.
Ever since my school days I had taken an interest
in historical reading, and I now planned a study
of the most important historical works, from Herodotus and Thucydides
to Macaulay and Ranke. For a moment I also thought
of answering a prize question the University had set
on a subject from the Danish Middle Ages; but nothing came
of it. Philosophy was pressing, and as a contrast
to all the theology I had worked through, I felt
a need to immerse myself in Greek philosophy. As I wrote some
% TR: "Doktorvita" = the curriculum vitae printed with the doctoral dissertation; rendered "doctoral vita".
years later in my doctoral vita: ``Classical antiquity
stood out for me in a fresher and clearer light precisely
through the contrast with the scholastic period I seemed
to myself to have lived through; I took real delight in
reading its historians, its poets, and soon especially
its philosophers as well. I thereby gained a deep conviction
of the independent significance of science and of humane
life.'' Here too Kierkegaard's guidance was constantly at work.
He had clearly set forth the Socratic-Platonic
conception of thought and life as a power of the spirit which one
had to come to terms with before one could, with full right and with
clear consciousness, arrive at a stage that really
was higher. I now read the Greek philosophical works
as a connected whole, and in the following years this became my main study.
In my teaching I exchanged, some years
later, history and Danish for Latin (in the middle forms)
and Greek (in the top forms), and so a certain
connection came about between my life at school and in the study.
Besides at ``Maribos Skole'' I also taught
at the ``Lyceum'' and for a short time at ``Bohrs Skole''. I was
% --- p. 57 ---
\opage{57}very fond of teaching, especially in the top forms, and
there I won friends who even after their school years have
shown me faithful devotion.

\whiteline

In the autumn of 1865 the young girl Emma Pape, who
a year before had gone away as my ``sister'', came to
town again to spend the winter there. As soon as I heard
that she was coming, it was clear to me that a decision must now
come, since she was in reality something other
and more to me than a sister. I was afraid of
a decision, not only as every young man doubtless is
% TR: "Enten-Eller" (the choice, not the book) rendered "Either/Or", keeping the echo of Kierkegaard's title.
when he sees an Either/Or approaching, but also for
reasons bound up with my religious stage. I
% TR: "Indrømmelse" = Kierkegaard's demanded "admission" (that official preaching is not New Testament Christianity); rendered ``admission''.
had for my part made the ``admission''
Kierkegaard demanded, and which prevented me from becoming
a pastor. But this admission did not exclude unrest
in the face of the question whether full participation in the
interests and tasks of natural human life was compatible
with a strictly Christian conduct of life. The strong passages
in the New Testament concerning marriage, and against marriage,
I could not get away from. And there now arose, called forth
by the Either/Or that was approaching, a last
struggle within me. The seventh chapter of First Corinthians
% TR: "Kierkegaards „Øjeblikke“" (plural: the issues of Øjeblikket) rendered "Kierkegaard's issues of ``The Moment''" here and below.
confronted me in all its might, and Kierkegaard's issues of ``The Moment''
lay on my desk. I was not ``seeking'' in the same
sense as Kierkegaard's ``cand. theol. Ludvig From'',
but I sought quite personally to make clear to myself what the consequence
was of what I held fast to as the highest.
It was the most decisive turning point in my life that I
stood at. The natural, human need prevailed, and
with that the foundation was in reality laid for my whole
later way of taking life. As I had earlier
flung the theological commentaries across the floor, so
I now flung Kierkegaard's issues of ``The Moment''. Those too I picked
up again, but from now on my occupation with
Kierkegaard was directed at finding the psychological and historical
% --- p. 58 ---
\opage{58}presuppositions of his appearance on the scene, and at gaining purely
human strength for my life by drawing the consequence
of the fact that subjectivity is the truth. ---

And then one afternoon in the autumn, while
I stood bent over my Thucydides, it came to stand
before me: now or never. Without giving myself time to
close the book I hurried out to my uncle's country house at
Frederiksberg, and there, on the starlit October evening,
the pact was made that began the short domestic
% TR: the square brackets are the author's own, as printed.
happiness that was to fall to my lot [in my young
days].

My fiancée stayed in Copenhagen over the winter. In the
following five years that passed before we could marry,
she was in Copenhagen roughly every other year, now in my
parents' house; the other years I spent my holidays
% TR: "Provst" = rural dean, rendered "Dean"; "Provstinden" below = "the dean's wife".
either at her uncle's, Dean Wulff's, parsonage in
Gunderup near Aalborg, or with her parents on
their farm in Vendsyssel.

Only on reading her letters from the time of our engagement again, many years after her death,
has it really dawned on
me how unfortunately placed she was in those years when she
was not in Copenhagen. Her uncle and aunt in Gunderup
were good and well-meaning people. But the dean's wife
was of the old school. The whole day was taken up with
housework and needlework. Walks were a luxury, and
so was reading or writing letters. At home with her
parents Emma had more freedom to write, read and
go for walks; but her mother was very difficult to get along with,
and there were many painful times. Several times
% TR: "Plads som ung Pige" = a living-in post as a young woman of good family helping in another household; rendered "a position as a young lady in a stranger's household".
Emma wished to take a position as a young lady in a stranger's household,
but I was always against it. At the time I did not really
make vivid to myself how difficult her situation was.
Good times we did have, though, partly in Copenhagen, partly in summer holidays
at her parents' farm in a characteristic
Vendsyssel landscape, with a wide view over the rugged and
hilly surroundings.

% --- p. 59 ---
\opage{59}In the winter of 1865---66 a discussion was carried on in the press
about Rasmus Nielsen's doctrine of faith and knowledge as two entirely
different principles which, precisely because of their difference,
were supposed to be capable of being united practically and personally.
I had been led to share this view by
my own personal development in religious respects, as well
as by my philosophical studies. The latter
needs a brief explanation.

In this country Hegel's philosophy was still regarded, around 1860,
as a high point of thought. There
a grand attempt had been made to show that, just as in the human
spirit one thought is understood only through its relation to
other thoughts, so in nature, in history, in the world of
experience altogether, no phenomenon, no event is
isolated, but everything is connected with everything else in the most intimate way.
As one thought proceeds from another, so
one phenomenon of experience proceeds from another. And Hegel
believed that it was possible to develop the mutual connection
between thoughts in such a way that it would at the same time come
to express the connection between things. From this
sprang his great confidence in pure speculation. Human
thought could not be something standing alone
in the world; its laws had to be the laws of the world.

In this country this doctrine had been the object of searching
criticism by Sibbern and Kierkegaard. The hints
these thinkers gave of a more critical and empirical
philosophy went unheeded. For the time being,
then, the situation was that both science and religion,
each on its own, laid claim to giving an absolute solution
of the riddles of existence. Here, for Hegel, the matter stood thus:
% TR: "Forestillinger" = Hegel's Vorstellungen, rendered "representations".
that the religious representations contained, in the form of the image,
the same thing that philosophy set forth in the pure form
of thought. Trusting in this view, many
Hegelians believed themselves to be in good agreement with the doctrine of the Church,
% TR: printer's error "Rasmus Nielsens hørte" (genitive) rendered by its sense, "Rasmus Nielsen".
and Rasmus Nielsen for a time belonged to their number.
But a great reversal came about in his view under the
% --- p. 60 ---
\opage{60}influence of Kierkegaard's writings. He broke with his former
associates, notably Martensen, and now asserted
the absolute independence of religion over against science.
They now came to stand as two different worlds, and
the human being had to live in both of them. This view
was attacked from two sides. Theology continued to demand
to be science and saw a danger in religion's being relegated
to the world of the purely personal and subjective. From
the opposite side it was above all Georg Brandes who came forward, in various
newspaper articles and later in a small work ``Dualismen i vor
nyeste Filosofi''. I took part in the controversy in a very clumsily
shaped work, ``Filosofi og Teologi'' (1866), which was in
the main a report of the discussion that, after
% TR: "Martensens Dogmatik" (Den christelige Dogmatik, 1849; unquoted in the print) rendered "Martensen's Dogmatics".
the appearance of Martensen's Dogmatics (1849), had been
carried on about the relation between philosophy and theology. My relation
to the whole matter I expressed some years later in my
% TR: printer's error "dennne" rendered by its sense, "this".
doctoral vita thus: ``I have always taken this view
[Rasmus Nielsen's conception of faith and knowledge]
more from its skeptical and negative side, more as a
means to a personal solution of the difficulties, than
as a positive doctrine of its own.'' That is to say, I was more
a Kierkegaardian than a Nielsenian. Brøchner later
expressly drew attention to the opposition between
Kierkegaard's and Nielsen's standpoints. And this opposition
% TR: "( som" (stray space after the parenthesis) rendered "(as".
became (as I have shown in ``Danske Filosofer'')
more and more pronounced in the course of Nielsen's continued development
of his thoughts, until at last (in his work ``Om
Aandsdannelse'' 1877) he entirely disowned Kierkegaard and
at bottom conceded that the latter's theological opponents were right.

For my part the situation was that I
was on the way out of theology, and that the question now was
what conception of science and what personal view of life
I could arrive at. A new period of apprenticeship now began
for me. The task was to gain a foundation for my further development
in both scientific and purely personal respects.
% TR: p. 60 ends mid-sentence, "Dette var jeg mig naturligvis ikke bevidst dengang, [p. 61 ...]"; rendered "Of this I was naturally not conscious at the time," and left open for p. 61; p. 61 continues it.
Of this I was naturally not conscious at the time,
% --- p. 61 ---
% TR: p. 61 opens mid-sentence; the sentence straddling pp. 60--61 reads
% "Dette var jeg mig naturligvis ikke bevidst dengang, men naar jeg ser
% tilbage finder jeg en uvilkaarlig Stræben efter et Fodfæste."  The English
% begins at "men" (but).
\opage{61}but when I look back I find an involuntary striving
for a foothold. About eight years passed before such a
foothold was reached. I had to come to know the distinctive
character of the various main problems and thereby
understand both their connection and their independence.
A short draft for a course of lectures in the autumn of
% TR: "Indledning til Filosofien" (Introduction to Philosophy), a lecture
% title, kept in Danish like the other titles.
1874 on ``Indledning til Filosofien'' contained the basic features
of everything I have since set forth at greater or lesser
length.

But I am getting ahead of myself. --- At the point I have now reached,
my studies of Greek philosophy were gathering more and
more around a definite task, which I treated in a
work, ``Den antike Opfattelse af den menneskelige Villie,''
which I defended for the philosophical doctorate in 1870.
Some years before, I had answered the University's philosophical
prize question (for 1867---68) and won the gold medal.
This question, too, concerned the human will,
since it called for an examination of the dispute over the freedom
of the will that was once (1824 ff.) carried on in our
literature between the forensic physician Howitz, Sibbern, A. S. Ørsted
and others. I am not satisfied with the solution I hinted at,
and in the judgment of the examiners (Nielsen
and Brøchner) the essay had its essential merit in its thorough
characterization of the various positions that
had made themselves felt, of which Howitz's and Sibbern's
had particularly interested me.

From this time on I came into a closer relation with Brøchner.
When I had won the gold medal, he advised me
to travel abroad, and my father then made me a present of a
journey. I wavered somewhat as to where I should
go. Alongside my other work I had tried
to find my bearings in the development of philosophy after
Hegel's time in Germany, and here the Göttingen philosopher
Hermann Lotze had particularly attracted my attention. I
% TR: Danish "tænkte" is partly supplied by the editor (\supplied{tæn}kte,
% scan loss); translated as the supplied word reads ("tænkte mig" = considered).
considered the possibility of studying with him in Göttingen.
But I judged that I could come to know this subtle thinker and
% --- p. 62 ---
\opage{62}brilliant author just as well through his books as through his lectures, and that a stay in a small German university town would not be able to enrich me in any other field than the purely scientific. I thought that I would get a more many-sided yield from spending a winter in Paris, where one could not only study, but where art and literature and the pulsing life of a capital offered opportunities for a more many-sided education. I spent a very full winter in Paris (1868---69), attended lectures (among others, Taine's), worked on my dissertation, was a constant visitor to museums and theaters, and used my eyes to observe life in this center of the world. Here I came to know other currents of thought than those that stemmed from Germany and had so far had the greatest influence on those at home who occupied themselves with philosophy. I became especially aware of the positivist school, represented partly by the great figure of Comte, who, beside and as the opposite of Hegel, was the greatest systematic thinker of the nineteenth century, partly by Herbert Spencer's bold but magnificent generalization of the concept of evolution. It was some time, however, before this school came to have a more decisive significance for me. I have always been slow to turn. Thoughts had to ripen quietly in me before I could work with them.

\whiteline

So far in this autobiography I have relied essentially on my memory. But to give a more immediate impression of the period of my life with which this section is concerned, one closer to the sources, I will insert some fragments of letters to my fiancée in the years 1866 to 1870, letters which I found among her things after her death, and which --- together with the letters from her that I had kept --- had now, in my seventy-fourth year, made it possible for me to put myself back into my inner and outer circumstances of those days. I pass over family matters and other things that are not of interest
% --- p. 63 ---
\opage{63}for the task I have set myself in this work. The passages from the letters are given in chronological order.

\whiteline

30 September 1866. --- If one tries to keep life going on idyllic, contemplative harmony, one will surely run up against reality and learn that the reconciliation was only in thought or in imagination \ldots{} That is why I have become a disciple of Rasmus Nielsen, because he so strongly stresses the opposition between faith and knowledge. They are, after all, two entirely different states of the soul. Faith is an ``unchanged mind,'' but knowledge belongs to the realm of contemplation.

% TR: 2 Tim. 3:12 is given in the King James wording, the Danish quoting its own Bible.
% TR: "lade, som om den hellige Grav var vel forvaret" -- proverbial Danish for
% complacent assurance that all is safe; kept literally ("the Holy Sepulchre").
25 November 1866. --- How gently and calmly life spreads out before us when we think of those to whose lot the great struggles and spiritual trials fell \ldots{} ``All that will live godly in Christ Jesus shall suffer persecution.'' (Second Epistle to Timothy) \ldots{} If the conditions of our life are to be the truth of our life, we must surely understand ourselves in the face of the great exemplars, and not, in obscurity, act as though the Holy Sepulchre were well guarded. What is required is humility and sincerity \ldots{} Then the raising up and the inward self-affirmation will not fail to come.

% TR: Danish prints "ngen" (printer's error); rendered by the evident sense "ingen" (no).
1866 (The date is missing). If I seldom go to church, it is partly because I know of no pastor who is sincere toward real Christianity. There are passages in the New Testament that one never hears any pastor speak about. It is by no means because I think I can be sufficient unto myself that I never go to church.

9 December 1866. --- You said recently that I seldom mentioned the family at home. That is not because I have not lived with them, and I mean not merely in outward contact but a real, inward life together. On the contrary, I have always felt a heartfelt joy in them, and I believe I understand them far better than before. A couple of years ago things were really rather cramped for me in that respect, as you surely know and probably saw yourself. I often felt so isolated and forsaken, so different from the others. But now I am certain that
% --- p. 64 ---
\opage{64}the fault must have lain with me. That there is a difference is certain enough; but what does even that really matter, as long as the right eye is there to see it? --- I could have become a bad sort if I had been allowed to go on the way I was a couple of years ago. But there was one who prevented it, who made my nature more open and gentle and drove away the coldness that for long periods could rest so heavily on me. It is to you I owe it that I can once more go about here at home with a free and glad mind among parents and brothers and sisters, for whom it must surely often have been a torment to look at such a mope as I was then. I do not at all think that I am quite discharged from the hospital --- I probably never shall be, and that is probably as it should be; but it is certain that I become afraid of myself when I think of how I was three years ago \ldots

This autumn I have come to know Marie (my sister). The truth probably is that I have needed her more because I did not have you. What a genuine, unaffected, childlike soul she is! There is gold in her, and it will stand its test. And the test may be hard; she knows only prosperity, after all --- but that, too, is a kind of test.

1 March 1867. --- You once said at Christmas that I never wrote anything about my studies \ldots{} So far as I can come to clarity with myself, I am convinced that it is the one right thing I am doing when I strive for a more comprehensive and deeper-reaching cognition and understanding in philosophical respects. It \textit{must} lead to a goal that will be of great importance for me personally; I feel so surely that my life would lack all firmness and would waver if I rejected these questions and left them standing in an indefinite fog. In that respect I harbor no doubt, although the goal naturally often seems to hide itself from me, and I seem to myself so powerless and so wretched. It is another question, on the other hand, whether I shall succeed in reaching a result such that it would also have some significance outwardly. One cannot, after all, go further than one's ability reaches. But do you not agree with me that the
% --- p. 65 ---
\opage{65}first consideration is the most important? (There followed a discussion of practical prospects; I have in mind to take the doctorate, work as a teacher at a school, and one day take over one of the private Latin schools here in town.) To become a university teacher is something that in our small country is, after all, a matter of sheer luck, quite apart from the demands it makes on gifts and personality.

7 April 1867. --- Now I am to lose Skovgaard. You may believe that I shall miss him greatly. He is to be private tutor with a Count Juel-Vind-Frijs on Lolland \ldots{} So I shall not hear his vigorous step on the stairs and loud thumping on the door, which for so many years now has announced a dear visit to me. It is no small benefit to have someone whom one can truly call one's friend in the full sense of the word, with whom one can speak from the heart and feel that one is wholly understood.

16 June 1867. --- You call me ``true.'' I think that this truth in me can turn into a timidity that perhaps slackens me not a little \ldots{} To be true is, to be sure, what I would wish above all. This truth, which is one with inwardness, is after all the only condition for our being able really to call anything our own \ldots{} Above all, untruth is so hard to get at when it hides behind an inspiring resolve.

8 September 1867. --- Sometimes I have a feeling that I am on a right way --- that is to say, on a way that can lead me to what I believe I ought most properly to attain.

15 September 1867. --- Daily life has its high points and its great falls as regards ideality. If only one could hold fast to the conviction that it is by the high points one is to judge, and that the radiance which falls from them over life is the true light in which life is to be seen.

22 September 1867. --- (About Julius Petersen). What I have often admired in him as energy and strength of will is, I believe, really more a restless striving, which can accomplish much
% --- p. 66 ---
\opage{66}when the goal seems to lie close at hand, but which collapses into itself when patient endurance is called for. (Out of disappointment at having been passed over in the appointment to a post as junior hospital physician, he thought of giving up the continuation of his studies and settling in the provinces.) It would be painful to me to see someone in whom I have had as much confidence as I have had in him give way to the pressure and obstacles of life.

13 October 1867. --- What I strive for, what more and more becomes for me a matter of spiritual welfare, is to work my way through to a firm view of the world. On that, of course, one must work throughout one's whole life, but I would still like in my youth to have carried through a scientific work such that I could have a good foundation in it \ldots{} That is what stands before me as the work of \textit{my} life.

27 October 1867. --- I always think that it is an idyll that is to be performed in life, and that is why things so often go so badly for me.

7 November 1867. --- It may perhaps interest you to hear in what direction my reading mainly runs. When the question of faith and knowledge, which tormented me when I was studying theology and which has since continually presented itself to me anew, led me to a closer occupation with the history of philosophy, I gradually became convinced that without having grasped Greek thought in its distinctive main forms one would not have the right idea of what it really is all about. For the Greeks had reached the highest perfection of spiritual development that could be reached under purely human conditions. With Christianity, faith came into the world, and new forces were set in motion, but everything we have in the way of human \textit{forms of culture,} all of it points back in its main features to the Greeks. I have therefore already been occupied for half a year with Plato and Aristotle, and if I succeed in carrying out my plan, I shall one day use a subject from this field for a doctoral dissertation. Alongside this
% --- p. 67 ---
\opage{67}I busy myself with history and Icelandic, which I need for my teaching.

14 November 1867. --- Lately I have read a good deal of Shakespeare aloud, especially on Sunday evenings, and also --- which is a great pleasure --- heard Poul Kierkegaard read. About a fortnight ago he, Troels Lund and I sat here one evening reading Shakespeare until the clock was almost one, without our noticing it.

24 November 1867. --- A change that has taken place in me this autumn consists in this, that I can no longer agree with Rasmus Nielsen, to whom, as you know, I had previously on the whole attached myself. Partly a wavering in himself, partly a telling critique by Professor Brøchner, partly and above all a renewed testing of the matter itself (not least in conversations with Poul Kierkegaard) have shown me the untenability of his view. Yet I have not thereby been brought closer to the ecclesiastical position \ldots{} On the contrary, it has at the same time dawned on me that on several points I had held opinions that had no true root in me but clung on as something external and handed down. Christianity will --- of this I am fully and firmly convinced --- always be to me the highest truth, the one in which alone I can find the riddles of life solved. But only what I can personally experience and appropriate can be truth for me; and I have come to the conviction that whatever one supposes oneself to possess that goes beyond this rests on unconscious hypocrisy. The view whose untenability I have now come to see has of course not been without significance; it was a stage of development that had to be gone through.

18 January 1868. --- For almost the whole of the new year I have been so listless and hypochondriac that I have often been at my wits' end \ldots{} It is never anything outside me that I need in those states; it is composure and clarity, and above all strength to keep myself under proper discipline, that are lacking.

(From the year 1868 there are not many letters, since we were
% --- p. 68 ---
\opage{68}together in Copenhagen from March on, until at the beginning of December I went abroad.)

19 June 1869. --- (About ``the period that lies before my engagement and even in part extended into it''). I have sometimes called this period my sentimental period, and with a certain right, but I believe it is more correctly designated by calling it the ascetic period. Happiness and joy then stood before me as unentitled, or at least suspect, guests at the table of life. Penance and renunciation were the leading thoughts on which the view of life rested. I do not say that at that time I lived in every respect as an ascetic; but the corresponding anxiousness was nonetheless present \ldots{} Had I lived in a Catholic country, I would, at least at certain times, not have been far from the cloister \ldots{} If I have become freer, more unprejudiced and more independent, I have also become more humble than I was then.

% TR: "en aandrig Livsfornødenhed" as printed; "aandrig" (normally witty,
% ingenious) is rendered in the sense the context requires ("of the spirit").
11 July 1869. --- This day has been good to me. I have truly experienced what a joy it is to be busy with a work that lays claim to one's interest. It is a necessity of life for the spirit. However modest what one accomplishes may be, there still come moments when one can have the consciousness: I too have felt the great thoughts on which existence rests and by which it is carried forward toward its goal. Even if I am not among those to whom it is given to grasp these thoughts and set them forth in a new and awakening form, my enthusiasm has been no less for that. Oh, if only one could once truly collect oneself, truly concentrate all one's powers and push aside every hindering power! Perhaps there might after all be a stone which it could be my lot to add to the great building.

18 August 1869. --- We have both, each in our own way, worked our way through to greater clarity and firmness \ldots{} I have often noticed with wonder how you saw things in the same way as I, and I am thinking especially of the freer religious views, which gain more and more significance and truth for me.

% --- p. 69 ---
\opage{69}26 June 1870. --- We now stand at the end of our engagement. How much has been experienced and lived through in that time! I really believe it is only now that I have become truly independent and thus a real man \ldots{} Independence consists in this, that things no longer unroll before thought like a panorama, but that one takes a hand oneself in order to draw forth what is true and eternal.

\whiteline

% TR: "Adjunktpladser": posts as Adjunkt, the lower rank of grammar-school master.
In accordance with my plan of going into school work, I applied repeatedly for posts as assistant master, and it was a disappointment to me that I did not succeed in getting such a position, now that I wished to set up house. Yet it would probably then have turned out that I would have had to give up continuing my studies on the scale I had envisaged. And after I had defended my dissertation for the doctorate in February 1870, I applied for and obtained the Smith scholarship, which I could expect to keep for some years if I continued my studies. On that, and on the income from my school teaching, together with an allowance from my father, I married in the autumn of 1870.

In the spring semester of 1871 I began to lecture as Privatdozent at the University. My endeavor to find my bearings within contemporary philosophy led to the publication of two books, ``Filosofien i Tyskland efter Hegel'' (1872) and ``Den engelske Filosofi i vore Dage'' (1874). Through a critical examination of English philosophy, especially as it was represented by Mill and Spencer, I was led to see the significance of a thought that had lain behind Kierkegaard's philosophizing and had been brought forward consciously still earlier by Kant, namely the thought of the unifying and unity-seeking character of the life of human consciousness. Later I took from here the hypothesis I made the foundation of my Psychology. In the autumn of 1874 I gave a course of lectures on ``Indledning til Filosofien,'' in which, from the standpoint thus gained, I characterized the various philosophical problems. From this time on it gradually \ldots
% TR: p. 69 ends mid-word with "Det blev fra denne Tid af efterhaan-" (efterhaanden,
% "gradually"); the word and sentence continue on p. 70, not yet transcribed.
\gap{pp.~70--76}
% [text to be added: pp. 70--76]
% p. 77 opens mid-word: "overnat-/tede" (spent the night); p. 76 not yet transcribed
% --- p. 77 ---
\ldots{} \opage{77}spent the night, so as to walk the next day to my parents-in-law's
farm, where I found wife and child in the best of health.

Every such walking tour I have undertaken has let me grow more intimately
into the nature of our country and feel myself one with it, on its sandy
coasts, in its lush forests and among its hills. Before and after the time I
% TR: Standard English exonyms (Copenhagen, Zealand, Funen, Jutland), as elsewhere in the repo; all other Danish place names as printed.
am speaking of here, I have walked through North, South and West Zealand,
Møen, Bornholm, southern Funen and great stretches of Jutland. For the
nature of Jutland in particular I have, as I have already said, a love, and a
summer vacation does not seem normal to me unless I have renewed my
acquaintance with it. Vendsyssel with its wide views, its hills, meadows and
streams, central Jutland with heaths and forests, heather and bogs, West
Jutland with dunes and the roar of the surf, East Jutland with the glories
of the Bay of Aarhus and Vejle Fjord --- all this not only gladdens me while
I am there, but images remembered from these regions often come back in my
fancies. In Steen Blicher's tales it is not least his splendid descriptions
of nature that have made him one of my favorite writers.

\whiteline

When, after this Jutland excursion, I return to my life in Copenhagen, there
is one figure who towers before me in my memory of the years around 1870. It
is Hans Brøchner. I had come into contact with him in the summer of 1867,
% TR: "Disputats" = the public defense of a doctoral dissertation; rendered "doctoral disputation" here and below.
when, after hearing me speak as opponent at a doctoral disputation, he sent
me a kind and encouraging greeting. From that time on I visited him often. I
got advice from him concerning my study of the history of philosophy, and
when I had written a critical review of one of his books, he sent me his
rejoinders and asked me to come one evening and discuss them with him. In
the following years, until his death (1875), I visited him regularly once a
month in his lonely home, where he lived as a widower with his
% --- p. 78 ---
\opage{78}two children. He associated with none of his contemporaries, except
% TR: "Geheimeraad" (Privy Councillor) is a Danish title of rank, left untranslated.
for the friend of his youth, Geheimeraad Vedel. Of the younger men only
Georg Brandes came to see him besides me, after his closest disciple, the
amiable and gifted Kristian Møller, had died. When one visited him, it was
as a rule not philosophical subjects that came up. In his solitude he felt a
desire to hear news from the academic world and from the world at large. I
can still see him sitting in his rocking chair or walking up and down the
floor with his characteristic long and determined strides. His unusually
handsome face could then light up with goodwill and interest. Often he told
% TR: Danish prints "sine unge, tidligt afdøde Hustru" (plural possessive with a singular noun); rendered by the evident sense, "sin".
of his travels or of his young wife, who had died early. Of Costanza, the
young Italian woman he had been in love with in his youth, he never spoke to
me; it was only when I published his letters that I came to know this
episode of his life. In the letters it appears, as he depicts it, as a
romance full of feeling. It seems to have been the intervention of a father
confessor that prevented him from bringing Costanza up to the North. It
would surely have been a hazardous venture, too. --- When one rose and was
about to take one's leave, he became especially lively and had many things
to ask about and talk over; it was as though he was loath to let go of the
few links that connected him with the world outside.

He was a philosophical hermit. The circle of thoughts he had made his own
was his spiritual home. He had no great gift for communicating himself in
free form, varying his expressions according to situations and
personalities, and he had once and for all summed up his ideas in fixed,
very abstract forms. But both in lectures and in conversation a tremor could
come into his voice and his expression, a sign of the deep feeling that lay
behind and bore witness to the connection of his thought with his whole
personality. He was a knight of thought, the last representative among us of
% TR: Danish prints "Tillidt" (hyphenated Til-/lidt); rendered by the evident sense, "Tillid" (confidence).
the idealism which, in the confidence that the innermost core of existence
could find, and had
% --- p. 79 ---
\opage{79}found, expression in the form of thought, also found in this form an
expression for the highest striving of personal life. For those who came
after him in the field of philosophy, the problems have become more complex
and manifold. In one of his last lectures Brøchner declared himself an
adherent of ``the idealist philosophy which Germany has the honor of having
% TR: Danish prints "denne Filosofis Hovedtænker" (this philosophy's principal thinker; the print, checked at 500 dpi, clearly has "æ"); rendered by the evident sense, "Hovedtanker" (principal ideas), a one-letter misprint.
produced,'' and there can hardly have been many in whom this philosophy's
principal ideas so pervaded and ennobled the personality as in him.
The image of ``the quiet professor,'' as he has been called, has followed
me on my way and reminded me what value the life of thought can have, even
beyond the purely intellectual domain, when there is earnestness in the work
of thought. His friendship was an honor and an encouragement to me, and it
was with deep emotion that I took leave of him on his deathbed. It was the
day before he died that I saw him for the last time. He spoke with complete
calm of his passing, and arranged with me how various matters were to be
settled. When he bade me farewell and thanked me for my friendship, I was
more overcome than he. His housekeeper related later that when I had gone he
had said quite calmly: ``Take that glass away, for Dr. Høffding's tears are
in it.''

A little myth took shape about his last moments after he had died. He died
on 17 December (1875) toward the end of the afternoon, that is, in the
darkest time of the year. Among the last things he said was: ``We must have
light in here.'' In the period just after his death I heard from several
quarters that he was supposed to have cried out for light, and there were
those who added that it was surely no great wonder that that man cried out
for light! --- The myth, however, seems to have died a natural death. For a
time it got some nourishment from what happened at his funeral. At that time
no purely civil funeral was possible. He had wished to be buried from the
University, and that a colleague should speak; but if that could not be
done, he had named a particular pastor, who had spoken over his young friend
Kristian Møller
% --- p. 80 ---
\opage{80}and had shown great tact on that occasion. But the pastor in question
seems to have misunderstood the situation and read more into the deceased's
highly conditional wish than was warranted. I thought of correcting the
misunderstanding, but gave it up when Brøchner's oldest friend, Geheimeraad
Vedel, firmly advised against it.

In the preface to my edition of the correspondence between Brøchner and
Molbech and in ``Danske Filosofer'' I have given a characterization of
Brøchner as a personality and as a thinker, and my work ``Den nyere
Filosofis Historie'' I have dedicated ``to the memory of my teacher and
friend.''

\secrule

If I were to try to describe my home life after I had set up house, I would
face the same difficulty I have spoken of earlier, as regards forming a
clear picture of those closest to one. The few years that this home life
lasted stand before me as an enclave in relation to the rest of my life. Not
as though it were a pure idyll. Looking back, I call myself to account for
many an hour when despondency and ill humor, which did not always have
deep-lying causes, were vented on the one I least of all wished to grieve.
But I was borne with, here as in my old home. And we developed side by side,
as had already begun in the years of our engagement. My wife was naturally
not at home in what occupied me. But when I managed to find an expression
for the human and personal side of the subjects my studies dealt with, it
was taken up with joy and understanding. I have always found it comical when
wives simply adopted their husbands' opinions, especially when, as is not
seldom seen, they come forward more dogmatically and didactically than their
male halves. And I may say that my wife had such a sense for what is genuine
and acquired by oneself, and such tact besides, that she went along only
with what she could make her own by the path of her own experience and
reflection. And in the land of poetry we could meet as in a common
% --- p. 81 ---
\opage{81}world. I read to her or with her Homer and Plato, Shakespeare and
Molière, besides the writers of the day, first and foremost Bjørnson and
Ibsen. I remember especially how taken she was with Homer. Many an evening,
when the question was what we should take up, she said: ``Let us have
something about the godlike heroes!'' Our household devotions now became
poetry. In the years of our engagement it had been the Bible in its
capacity as dogmatic authority. Hand in hand we had passed over from the
land of dogmas to the land of symbols and poetry. Now too we read the Bible,
but as great poetry of life.

Glorious summer journeys we made together, to Møen, to Aarhus and the
Silkeborg country, or we met at her parents' farm in Vendsyssel, she
traveling on ahead while I would usually make a little excursion, partly on
foot, and end up with her and her family. I have mentioned the walk along
Jammerbugten. Another time I went to Skagen. That was in the summer of 1875.
Skagen had at that time only just been discovered by Drachmann. No hotels or
villas yet waged war on free nature. One drove by mail coach from
Frederikshavn, partly right along the water's edge, and from the post
station one had to wade through deep sand to the thatched inn. Here I stayed
% TR: "Indremissionær": an adherent of Indre Mission (the Inner Mission), the Danish pietist revival movement.
almost as the only guest; an Inner Mission man was staying there too. So I
could roam the world of sand and sea undisturbed. The return journey I made
by carriage to Frederikshavn, and from there I walked the glorious road along
the shore, with the Bangsbo hills on my right hand, toward Sæby. That bright
summer afternoon still stands clearly in my memory, and not least the moment
when, a little outside Sæby, I met my wife, who had driven out to meet me. ---

In Copenhagen we associated first and foremost with Julius Petersen and his
wife Oliva, née Munch. He was now practicing in Nørrebro and quickly built up
a large practice, not least in the poor quarters. He was restlessly active,
and was already beginning his studies in the history of medicine. When he
came home from sick calls in poor homes, he could be indignant at the luxury
he thought he saw in
% --- p. 82 ---
\opage{82}their very meager homes, and he could sit sunk in ill humor and
discontent. But then suddenly a quite different mood could come over him,
and often, when we thought the barometer stood low, an invitation might come
to join a little party. Against his restlessness and his many swings of mood,
Oliva's level-headedness, clear understanding and firm will were of good
use. As the flock of children quickly grew, there was plenty for her to do.
And she always kept her house running well, so that it even became a
gathering place for his friends. Here we met Kristian Arentzen, Viggo Hørup
and Holger Drachmann and their wives, and in many long speeches at the
festive supper table the many questions of the day were debated. --- Troels
Lund I had got to know in my student years. From philology he had gone over
to theology, and after his examination he occupied himself, like me, with
studies of Greek philosophy and took his doctorate (1871) with an
interesting dissertation on Socrates. Later he went over to historical
studies, occasioned by an appointment at the National Archives (the
Geheimearchiv, as it was then called). Troels Lund's fine, harmonious
personality made him a pleasure to be with, and his lively wife together
with him made their home a place where we liked to come. And in those days
one could also have the pleasure of being surprised by a visit; it was so
cozy when Lund and his wife came and settled in for the evening, making do
with whatever the house could offer. --- To Heegaard's home, too, many good
memories lead back. At my doctoral disputation he was the opponent, and from
that time on we came to be close to each other. After breaking with Rasmus
Nielsen he remained a seeker, absorbed in eager study of Kant and the more
recent English philosophers; he also occupied himself a great deal with
mathematics. As his financial circumstances were straitened, he had to take
on exhausting work on the side. Among other things he wrote anonymously a
series of novels, in which he was strongly influenced by Dickens. For some
years they found \ldots
\gap{pp.~83--89}
% [text to be added: pp. 83--89]
% p. 90 opens mid-sentence: "filosofisk set var meget overfladisk" (philosophically speaking was very superficial); p. 89 not yet transcribed
% --- p. 90 ---
\opage{90}\ldots{} philosophically speaking was very superficial, --- and that it had
more eye for what had to fall before criticism than
for what retains its significance, and for the new possibilities
of positive value. For my part, the
main thing was that I had enough to do in getting
firm ground under my feet in the field where I
personally was to work. By the middle of the '70s I had arrived
at what one might call a critical positivism (and what
I have somewhere, from another point of view, called a critical monism).
It set me ever new tasks, which I
had not been able to foresee, but which gave me
constant work. I had, as I have said in ``Tilbageblikket'',
% TR: ``Tilbageblikket'' (``The Retrospect'') kept as printed, as a title.
a thread to spin further, whether the current of the time
ran in the one direction or the other. The fluctuations
of literature, which modern literary historians believe
they can follow from decade to decade, so that each decade
gets its own characterization, I have not gone along with, and
I have not really been able to trace these rhythms in the young
generation I have had to do with either. But then I am probably not
fine enough in my senses either.

Georg Brandes and I had been opponents in the controversy
over faith and knowledge. Later we often met in the house of
Troels Lund's father. When he, together with his brother, had
founded the periodical ``det nittende Aarhundrede'', I
became a contributor and thereby came into a closer, and for me
very beneficial, connection with him.

He invited me one evening after he had married, and
he and his wife also visited my home. I have always
followed him with admiration in his works of literary criticism.
For me, as for so many others, he has opened
new poetic and literary worlds, and by his gift for penetrating
characterization and analysis he has thrown light on
so many of the great figures of literature. I felt myself
very different from him in natural temperament. I took everything
more heavily and more awkwardly, and I lacked the self-confidence
that is the condition for bold public appearance. What I
% --- p. 91 ---
\opage{91}have been able to accomplish is owing to long and unflagging work,
and my tasks lay further from the general public. Moreover,
I have on the whole held the view that what
is to be changed or removed in the world of thought or of society
% TR: "erkendes" (verb): "be recognized", in its ordinary sense.
can only be recognized once the right thing is found.
It is difficult to find the right thing, but when it is found,
it becomes clear what the false is. And the forms in which
men's need for support and footing in life finds rest
can be changed only when this need itself changes; but
such change takes place very slowly when it is to be
thorough. Here too it is the truth that shows what
is false, positivity that grounds negativity.
As Spinoza says: truth bears witness both to itself
and to untruth (verum est index sui et falsi). And
Comte maintained that one only displaces what one
replaces. This way of looking at things meant that I stood
differently from Brandes toward religious and moral polemic.
But his criticism has in many ways benefited my
activity, since in many it awakened a need for
further orientation in the field of thought. I have,
to be sure, in not a few cases seen how the view
came up that it was enough to have brisk and audacious opinions,
without one's needing to care that philosophy demanded
a searching investigation.

What I am not least grateful to Brandes for from the time I am speaking of here, the Seventies,
is the advice he gave
me in connection with the contributions I sent to his periodical.
Thus he wrote to me on 7 June 1875: ``Allow
me, not so much as editor as in the capacity of a sympathetic
fellow striver, to put a couple of questions to you.
Have you ever made the experiment of reading what you
give to the press aloud to some man or other who has not
been to a grammar school? Have you read it aloud to a lady
who was not especially trained in philosophy? Have you put a mark
by the places he or she did not understand, and tried to rewrite
them in other words? I dare say that you have not
% --- p. 92 ---
% TR: "Og derpaa anker særligt over": the print has no subject; "he" (Brandes) supplied.
\opage{92}made this experiment. But why not?'' And then he complains
in particular about a passage in an article I had submitted to
the periodical, and which was undeniably expressed in a very
unpopular way. I found the complaint just; unfortunately it probably
also hits much of what I have written since, although I
really did set myself to write as accessibly to all as
possible. But there is, after all, a difference between philosophy and literary history,
and besides it has been my lot to develop
slowly and to come late to clarity. I could have
waited to write until everything could be expressed clearly (if
not for those who have no grounding at all); that
might perhaps not have been to the detriment of literature. But
I have always learned much myself by trying to set forth
my thoughts and the knowledge I had so far acquired. And
I had in Brøchner an example of what it can lead to
to wait too long with writing; he often regretted that
he had chosen to wait with his writing until he had
thoroughly worked through all his subjects, and when that
time seemed to him to have come, he was nearly fifty,
and he had grown wholly into the abstract expressions
that had been handed down from the Hegelian philosophy, and
which in themselves could have done good service as short
formulas in his solitary thinking. --- I have, though,
often been surprised later at how wide the circles
are that many of my books have reached.

When in the year 1876 I had published my little book ``Om
Grundlaget for den humane Etik'', which I shall come back to in another
connection, Brandes wrote to me
% TR: the nested „forklarede“ inside Brandes's quotation rendered with single quotes.
(24 Oct. 1876): ``You are, though on another foundation, the
`transfigured' Brøchner, and it would certainly have been
a great joy to Brøchner if he had lived to see this book
of yours. .... I congratulate you also on the form of your
book. You have gained extraordinarily here in freedom and suppleness.
You once said a despondent word to me about dryness
in your language, where you were far too strict with
yourself .... In my opinion there is still too little
% --- p. 93 ---
\opage{93}temperament in your style; you do not yet write sufficiently
with your whole personality, the whole human being.
But you are on the way to doing so.'' --- For the rest, Brandes
recalls in this letter the inscription on the papers of the Committee of Public Safety:
``Liberté, Egalité, Fraternité, ou la mort!''
--- and he thinks that I have hardly ever been so
young that I could have said this. In that he is quite right.
Yet it grated a little in my ears when he added: ``You
must have been born to some degree old.'' Of himself
he then says: ``I am afraid that I was once very
young.'' ---

The relationship in which I had come to stand with Brandes was
a source of much joy and encouragement to me. After Brøchner's death
I felt lonely, and since Brandes combined with his great artistic
sense an understanding of philosophical inquiries,
his letters were always welcome to me. Some years later
he wrote to me from Berlin: ``I am very glad
to be on such friendly terms with you; there is nothing I
would rather do than live in brotherly understanding with all those of my generation
in the North who inquire honestly.'' What was expressed here
so warmly found a full echo on my side.
And this feeling did not vanish in me even when, a
good many years later, it was put to the test by a serious and
sharp opposition between our views on an essential
point.

% --- p. 94 ---
\chapter*{III}
\addcontentsline{toc}{chapter}{III. 1874--1883}
\markboth{III. 1874--1883}{III. 1874--1883}

% TR: the lecture title „Indledning til Filosofien“ (``Introduction to Philosophy'') kept as printed.
\opage{94}In the autumn of 1874 I gave a course of lectures on ``Indledning
til Filosofien'', in which I set forth in all brevity
the conception of the philosophical problems and of the relations
between them that had now become settled in me. And
after that my authorship proper began, as
I now set myself the task of treating the particular problems
each by itself. The short program I had given
in the lecture course mentioned naturally had to undergo
many changes in the working out of the details, and
there was a great deal of work to be done in the particular fields. I
can divide my work over the following 40
% TR: "Fra 1875 til 1887 drejede sig om" printed without a subject ("det"); rendered "The years from 1875 to 1887".
years into three periods. The years from 1875 to 1887 turned on psychology and ethics
in the first place, from 1887--1895 on the history of philosophy, and
% TR: "Erkendelsesteori" = epistemology (repo term Erkendelse = cognition; the compound given its standard name).
after 1895 it is philosophy of religion and epistemology that
have especially laid claim to me. As the position of philosophy
still stands, it is of some importance that these different
fields, which in so many ways overlap
one another, can be treated by the same person. It is not
certain that this will always be possible, and the problems of philosophy
would then present themselves in more difficult forms. Our work we
must at any time take up as the questions present themselves to
us; only we must not forget that questions as well as
answers can take on a different character at different times.
Yet I cannot conceive of philosophy without a striving
to gather human thought-work under certain
general points of view. And this striving is prepared within
the particular sciences, the more these themselves seek to become
% --- p. 95 ---
\opage{95}conscious of their guiding ideas. There is then always a task
for philosophy as a comparative science.

The first tasks that presented themselves to me were: to
write an ethics on purely human presuppositions,
and to work out a psychology by which the study of the
human life of the soul, which since Sibbern's time had not
been cultivated in this country, could be resumed. From
this time on I felt my activity more and more as a continuation
of Sibbern's, or rather, I wished that it could
be regarded so. Much as I owed Rasmus Nielsen
and Brøchner, I saw nevertheless that one had to go other
ways, ways which ``old Sibbern'' had traveled over many stretches,
although a romantic-speculative background had at
many points determined the light in which he saw
the questions.

In the autumn of 1876 I published a little work ``Om Grundlaget
for den humane Etik''. I sought here to show that there
is a purely human foundation for the conduct of life, which
is given with the possibility that man can submit
to a law that points beyond the particular moments and
beyond the particular individuals. Although I believe that I got
some way further in clarity with that work, and that in particular the analysis
of the concept of authority may still have some interest,
there remained an indeterminacy, which was connected
with the fact that I did not yet distinguish between the motive that leads
us to call actions good or evil, and which I
% TR: Vurderingsmotivet / Handlingsmotivet = the motive of valuation / the motive of action (Høffding's terms of art).
have later called the motive of valuation, and the motive that in
the particular special cases makes us satisfy the demands of the motive of valuation,
and which I have later called the motive of action.
When I got full clarity about this distinction,
which is so often overlooked, I could proceed to the working out of
a full-length ethics. But that happened only a good many years later.

For the time being I became absorbed in psychological studies and treated
each year, from the autumn of 1876 on, some section or other
of psychology in my lectures as Privatdozent.
Gradually a considerable circle of listeners gathered around \ldots
% p. 95 ends mid-sentence; p. 96 not yet transcribed
\gap{pp.~96--120}
% [text to be added: pp. 96--120]

% ---- Chapter IV.  Indhold: "IV. 1883—1894 .... 121" (so read at the
\chapter*{IV}
\addcontentsline{toc}{chapter}{IV. 1883--1894}
\markboth{IV. 1883--1894}{IV. 1883--1894}
\gap{Chapter IV, pp.~121--174,}
% [text to be added: pp. 121--174]

% ---- Chapter V.  Indhold: "V. 1894—1904 .... 175".  pp. 175--223 (PDF 177--225).
\chapter*{V}
\addcontentsline{toc}{chapter}{V. 1894--1904}
\markboth{V. 1894--1904}{V. 1894--1904}
\gap{Chapter V, pp.~175--223,}
% [text to be added: pp. 175--223]

% ---- Chapter VI.  Indhold: "VI· 1904—1914 .... 224".  pp. 224--265 (PDF 226--267).
% TR: the Danish head on p. 224 prints "V" (the I did not ink); rendered "VI".
\chapter*{VI}
\addcontentsline{toc}{chapter}{VI. 1904--1914}
\markboth{VI. 1904--1914}{VI. 1904--1914}
\gap{Chapter VI, pp.~224--265,}
% [text to be added: pp. 224--265]

% ---- Chapter VII.  Indhold: "VII. 1914—1922 .... 266".  pp. 266--310 (PDF 268--312).
\chapter*{VII}
\addcontentsline{toc}{chapter}{VII. 1914--1922}
\markboth{VII. 1914--1922}{VII. 1914--1922}
\gap{Chapter VII, pp.~266--310,}
% [text to be added: pp. 266--310]

% ---- Chapter VIII.  Indhold: "VIII. 1922—1924 .... 311".  pp. 311--324 (PDF 313--326).
\chapter*{VIII}
\addcontentsline{toc}{chapter}{VIII. 1922--1924}
\markboth{VIII. 1922--1924}{VIII. 1922--1924}
\gap{Chapter VIII, pp.~311--324,}
% [text to be added: pp. 311--324]

\backmatter

% ---- CONTENTS [Indhold], p. [325], unnumbered.
\chapter*{\emph{CONTENTS}}
\addcontentsline{toc}{chapter}{Contents}
\markboth{Contents}{Contents}
\opage{325}%
\begin{center}
\begin{minipage}{0.85\textwidth}
\small
\hfill{\footnotesize Page}\par
\noindent\makebox[3em][r]{I.}\enspace 1843---1861\dotfill 7\par
\noindent\makebox[3em][r]{II.}\enspace 1861---1874\dotfill 38\par
\noindent\makebox[3em][r]{III.}\enspace 1874---1883\dotfill 94\par
\noindent\makebox[3em][r]{IV.}\enspace 1883---1894\dotfill 121\par
\noindent\makebox[3em][r]{V.}\enspace 1894---1904\dotfill 175\par
% TR: the Danish prints "VI·" with a raised point; rendered "VI.".
\noindent\makebox[3em][r]{VI.}\enspace 1904---1914\dotfill 224\par
\noindent\makebox[3em][r]{VII.}\enspace 1914---1922\dotfill 266\par
\noindent\makebox[3em][r]{VIII.}\enspace 1922---1924\dotfill 311\par
\end{minipage}
\end{center}

\begin{small}
% TR: "Paranteser" (so printed) = parentheses.
Sections I.---VI. were written down in the winter of 1915---16. --- Section
VII. was written in the year 1922. --- Section VIII. in the spring of 1924.
--- Before printing, a read-through led to the addition of some parentheses
and notes.
\end{small}

\secrule

\end{document}
